\chapter{Web I: C\texttt{++}23 to WebAssembly}
\label{ch:web-compiling}
\label{ch:web-emscripten}
\index{Emscripten}\index{WebAssembly}\index{web build}

CNA does not have a separate JavaScript implementation. Its Web builds send the same
C\texttt{++}23 framework through Emscripten and link it with SDL's browser port. Configuration
admits 24 renderer identities on Emscripten, but five are the web-focused evidence set examined
here: \texttt{WEBGL1}, \texttt{WEBGL2}, \texttt{CANVAS}, \texttt{HTML\_DOM}, and
\texttt{SVG\_DOM}. One configure still packages exactly one selected identity into a WebAssembly
application. That compact description hides three independent questions: can the source be
translated, can the program survive the browser main-loop hand-off, and has the selected renderer
actually been exercised? This chapter answers only the first. Chapters~\ref{ch:web-main-loop} and
\ref{ch:web-renderer-evidence} answer the other two.

\section{The toolchain changes the execution contract}

The root build selects \texttt{WEBGL2} by default when \texttt{EMSCRIPTEN} is true.  This is a
public renderer identity backed by the EasyGL family, just as \texttt{WEBGL1} selects another
configuration of that family.  \texttt{HTML\_DOM}, \texttt{SVG\_DOM}, and \texttt{CANVAS}
select distinct implementations.  The choice is still exclusive: one configure produces one
renderer, not a universal Web binary containing all five.

Exception handling is the first project-wide adjustment.  CNA uses C\texttt{++} exceptions in
ordinary control and error paths, so the Web build enables real WebAssembly exceptions with
\texttt{-fwasm-exceptions}.  This is not merely a compiler-acceptance flag.  It affects what
the generated cleanup machinery catches, and therefore participates directly in the
main-loop lifetime defect studied in Chapter~\ref{ch:web-main-loop}.

The build remains single-threaded.  No Web target enables Emscripten pthread support, and the
tree has no corresponding cross-origin isolation setup.  One CnaTests configuration uses
\texttt{ASYNCIFY=1} for asynchronous WebSocket/ENet behavior; that transformation is not
evidence of a worker-thread architecture.  A port that needs concurrent game, networking, or
asset work would have to define that architecture and its hosting headers explicitly.

\section{Assets must cross the virtual-filesystem boundary}

A native executable can discover files already present beside its working directory.  A
browser cannot.  Emscripten exposes a virtual filesystem, and a file exists there only if the
page creates it or the linker packages it.

At the audited revision, exactly two build sites use \texttt{--preload-file}: the 2D graphics
demo and the sound demo.  There is no project-wide asset manifest, no \texttt{--embed-file}
policy, and no custom shell file that silently fills the gap.  Consequently, successful
translation of a target proves nothing about the availability of its content at run time.
Each Web application must account for every filename that can reach
\texttt{TitleContainer} or \texttt{ContentManager}.

The same distinction applies to writes.  SDL's preference path resolves inside Emscripten's
in-memory filesystem, under \texttt{/libsdl}.  CNA does not mount IDBFS, call
\texttt{FS.syncfs}, or bridge its storage abstractions to \texttt{localStorage}.  Saved data
therefore survives within the current page instance but not a reload.  \texttt{StorageDevice}
and GamerServices persistence on a native platform must be described as volatile on Web until
a durable browser store is deliberately connected.

\section{Networking is a different topology, not a transparent port}

Emscripten's socket shim does not report a usable ephemeral bound port, so CNA's ENet host path
requests fixed port 61191 instead of \texttt{ENET\_PORT\_ANY}. That host role is meaningful for a
Node relay, not an ordinary browser tab: the browser-side client path replaces any listening host
with an outbound-only ENet client before connecting. The entire UDP-broadcast discovery
implementation is excluded under \texttt{\_\_EMSCRIPTEN\_\_}. A Web build can therefore preserve
SystemLink message transport without preserving native hosting and LAN-discovery topology.

The CnaTests-only \texttt{ASYNCIFY=1} setting lets synchronous-looking polling yield to the JS
event loop during WebSocket/ENet tests. It does not add listening sockets, make discovery
available, or propagate to ordinary demos. Build evidence for the network module must therefore
name the role and host, not merely say that ENet linked.

\section{What a successful build establishes}

A completed Emscripten link is meaningful evidence.  It establishes that the selected source
closure is accepted by the Web toolchain, that its C\texttt{++}23 surface is compatible with
the chosen SDK, and that the requested renderer's symbols can be resolved.  It does not
establish any of the following:

\begin{itemize}
  \item that all runtime assets were packaged;
  \item that save data survives a page reload;
  \item that a frame was reached after Emscripten took control;
  \item that generated WebGL shaders compiled in a browser;
  \item that the page displayed the selected renderer rather than merely loading its module.
\end{itemize}

This boundary matters because a bare-Node launch is a poor substitute for a browser.  DOM
renderers need \texttt{window} and \texttt{document}; SDL-backed renderers need browser host
services; WebGL needs a real graphics context.  Errors such as \texttt{window is not defined}
or an SDL initialization failure say that the wrong host was used, not that the renderer's
drawing semantics were tested.

\section{A reproducible build record}

Evidence for a Web build should record at least the CNA commit, emsdk version, renderer
identity, configure options, named target, and produced \texttt{.html}, \texttt{.js}, and
\texttt{.wasm} artifacts.  It should also list every preload mapping.  Without the mapping,
two apparently identical builds may have different runtime content.

For testing, the record should say whether the target was merely linked, launched under Node,
or served over HTTP to a named browser.  These are ascending evidence levels, not synonyms.
The current HTML-DOM workflow, for example, pins emsdk~6.0.3 and serves six pages to Playwright's
Chromium.  That is a browser execution result; a successful build of the other identities is
not equivalent to it.

\section{The portability lesson}

WebAssembly preserves much of the language and loses much of the process model.  The hard
porting work is therefore at the boundaries: ownership across a non-returning host call,
assets crossing into a virtual filesystem, persistence crossing into a browser store, and a
verdict crossing out of a process whose exit status the harness cannot observe.  Treating
those boundaries as explicit interfaces keeps ``Web build'' from becoming a claim broader
than the evidence behind it.
